\begin{tabularx}{\textwidth}{>{\centering\arraybackslash}X >{\centering\arraybackslash}X}
\textbf{TRUNG TÂM TMATH EDU} &
\textbf{KỲ THI KIỂM TRA CUỐI KÌ I} \\

\textbf{NĂM HỌC 2025 - 2026} &
\textbf{Dành cho lớp 9C5} \\

\fbox{\begin{minipage}{2.8cm}
\centering
\textbf{ĐỀ SỐ 1}
\end{minipage}} &
\textbf{Thời gian làm bài: 90 phút} \\
\end{tabularx}

\vspace{0.3cm}
\hrule
\vspace{0.6cm}

\renewcommand{\arraystretch}{1.2}
\begin{tabularx}{\textwidth}{|c|X|c|c|c|}
\hline
\textbf{Bài} & \textbf{Tên bài} & \textbf{File đầu vào} & \textbf{File đầu ra} & \textbf{Điểm} \\
\hline
A & Ma Trận Đề Thi Lượng Tử & \texttt{A.INP} & \texttt{A.OUT} & 6 \\
\hline
B & Ba Quầy Hàng & \texttt{B.INP} & \texttt{B.OUT} & 5 \\
\hline
C & Đánh giá nhóm gia đình & \texttt{C.INP} & \texttt{C.OUT} & 5 \\
\hline
D & Trạm Năng Lượng Tuyến Tính & \texttt{D.INP} & \texttt{D.OUT} & 4 \\
\hline
\end{tabularx}

\section*{Bài 1: Ma Trận Đề Thi Lượng Tử}

Trong năm 2145, Liên Minh Thiên Hà tổ chức một kỳ \emph{Đánh giá Năng lực Trí tuệ} cho các phi hành gia tập sự.
Để tránh gian lận thông tin giữa các thí sinh ngồi gần nhau, hệ thống trung tâm đã chuẩn bị nhiều \emph{ma trận đề thi lượng tử} khác nhau.

Các ma trận này được đánh số từ $0$ đến $p-1$, với $p$ là số lượng ma trận mà hệ thống kích hoạt trong kỳ thi đó.

Mỗi thí sinh được gán một \emph{mã định danh sinh học} là một số nguyên dương \textbf{không trùng nhau}.  
Hệ thống phân phối đề thi hoạt động theo quy tắc cố định:

\begin{quote}
    Nếu một thí sinh có mã định danh là $r$, thì ma trận đề thi được kích hoạt cho thí sinh đó có chỉ số
    \[
        (r - 1) \bmod p.
    \]
\end{quote}

Hai phi hành gia, Ram và Shyam, được bố trí trong cùng một khoang thi và hy vọng có thể \emph{chia sẻ dữ liệu thần kinh}.  
Điều này \textbf{chỉ có thể xảy ra} nếu cả hai nhận được \emph{cùng một ma trận đề thi lượng tử}.

Bạn được cung cấp mã định danh sinh học của Ram là $A$ và của Shyam là $B$.

Nhiệm vụ của bạn là xác định:
\begin{itemize}
    \item Có bao nhiêu giá trị của $p$ để Ram và Shyam nhận cùng một ma trận đề thi;
    \item Hoặc xác định rằng có \textbf{vô hạn} giá trị $p$ như vậy.
\end{itemize}

\subsection*{Dữ liệu vào}

\begin{itemize}
    \item Dòng đầu tiên chứa số nguyên $T$ $(1 \le T \le 100)$ — số lượng kịch bản kiểm tra.
    \item $T$ dòng tiếp theo, mỗi dòng chứa hai số nguyên $A, B$ $(1 \le A, B \le 10^8)$ — mã định danh của Ram và Shyam.
\end{itemize}

\subsection*{Dữ liệu ra}

Với mỗi kịch bản:
\begin{itemize}
    \item In ra số lượng giá trị $p$ sao cho Ram và Shyam nhận cùng một ma trận đề thi;
    \item Nếu số lượng đó là vô hạn, in ra \texttt{-1}.
\end{itemize}

\subsection*{Ví dụ}

\noindent
\begin{minipage}[t]{0.48\textwidth}
\textbf{Input}
\begin{verbatim}
1
2 6
\end{verbatim}
\end{minipage}
\hfill
\begin{minipage}[t]{0.48\textwidth}
\textbf{Output}
\begin{verbatim}
3
\end{verbatim}
\end{minipage}

\bigskip

\subsection*{Giải thích}

Với $A = 2$ và $B = 6$, hai phi hành gia sẽ nhận cùng một ma trận đề thi nếu các giá trị $p$ thỏa mãn là:
\[
    p = 1,\; 2,\; 4
\]

Do đó, có tổng cộng \textbf{3} giá trị như vậy.


\bigskip

\section*{Bài 2: Ba Quầy Hàng}

Trong một khu phố nhỏ, có ba quầy hàng quen thuộc mà mỗi gia đình thường ghé qua:

\begin{itemize}
    \item Quầy thực phẩm khô, nơi mỗi món có mức giá cố định.
    \item Quầy đồ uống, với nhiều loại chai khác nhau.
    \item Quầy trái cây, bày bán các loại quả theo từng mức giá.
\end{itemize}

Mỗi quầy đều niêm yết sẵn danh sách giá cho các mặt hàng của mình.  
Một khách hàng khi đi chợ luôn chọn \textbf{đúng một món} từ \textbf{mỗi quầy}.

Trong ngày, có nhiều khách hàng khác nhau đưa ra câu hỏi:

\begin{quote}
    \emph{“Với số tiền này, liệu tôi có thể mua đúng một món ở mỗi quầy sao cho tổng tiền phải trả đúng bằng số tiền tôi đang có không?”}
\end{quote}

Nhiệm vụ của bạn là trả lời từng câu hỏi đó.

\subsection*{Dữ liệu vào}

\begin{itemize}
    \item Dòng đầu tiên chứa số nguyên dương $n$ $(1 \le n \le 100)$.
    \item Dòng thứ hai chứa $n$ số nguyên không âm $a_1, a_2, \dots, a_n$ $(0 \le a_i \le 10^8)$.
    \item Dòng thứ ba chứa số nguyên dương $m$ $(1 \le m \le 100)$.
    \item Dòng thứ tư chứa $m$ số nguyên không âm $b_1, b_2, \dots, b_m$ $(0 \le b_i \le 10^8)$.
    \item Dòng thứ năm chứa số nguyên dương $l$ $(1 \le l \le 100)$.
    \item Dòng thứ sáu chứa $l$ số nguyên không âm $c_1, c_2, \dots, c_l$ $(0 \le c_i \le 10^8)$.
    \item Dòng thứ bảy chứa số nguyên dương $Q$ $(1 \le Q \le 2 \times 10^5)$.
    \item Dòng thứ tám chứa $Q$ số nguyên không âm $x_1, x_2, \dots, x_Q$ $(0 \le x_i \le 3 \times 10^8)$.
\end{itemize}

\subsection*{Dữ liệu ra}

Gồm $Q$ dòng.  
Dòng thứ $i$ in ra:

\begin{itemize}
    \item \texttt{Yes} nếu khách hàng thứ $i$ có thể chọn đúng một món ở mỗi quầy sao cho tổng tiền đúng bằng $x_i$;
    \item \texttt{No} trong trường hợp ngược lại.
\end{itemize}

\subsection*{Ví dụ}

\noindent
\begin{minipage}[t]{0.48\textwidth}
\textbf{Input}
\begin{verbatim}
3
1 2 3
2
2 4
6
1 2 4 8 16 32
4
1 5 10 50
\end{verbatim}
\end{minipage}
\hfill
\begin{minipage}[t]{0.48\textwidth}
\textbf{Output}
\begin{verbatim}
No
Yes
Yes
No
\end{verbatim}
\end{minipage}

\subsection*{Giải thích}

\begin{itemize}
    \item Không có cách mua nào cho tổng tiền bằng $1$.
    \item Có thể chọn các món $1+2+2$ với tổng giá bằng $5$.
    \item Có thể chọn các món $2+4+4$ với tổng giá bằng $10$.
    \item Không có cách chọn nào cho tổng tiền bằng $50$.
\end{itemize}

\section*{Bài 3: Đánh giá nhóm gia đình}

Trong một khu dân cư, có $N$ gia đình sinh sống dọc theo một con phố.
Mỗi gia đình được ghi nhận một \textbf{chỉ số điều kiện sinh hoạt} (có thể âm hoặc dương),
lần lượt là: $A_1, A_2, \ldots, A_N$

Ban quản lý khu phố thường xuyên tổ chức các \textbf{buổi sinh hoạt nhóm},
mỗi nhóm gồm \textbf{ba gia đình khác nhau}, được chọn sao cho thứ tự xuất hiện
của họ trên con phố là tăng dần.

Với mỗi nhóm gồm ba gia đình, ban quản lý đánh giá mức độ ổn định của nhóm đó
bằng cách lấy \textbf{giá trị thấp nhất} trong ba chỉ số sinh hoạt của nhóm.

Tập hợp toàn bộ các mức đánh giá thu được từ mọi nhóm ba gia đình
được gọi là \textbf{danh sách đánh giá nhóm}.

Do nhu cầu thống kê, ban quản lý đưa ra $Q$ yêu cầu độc lập.
Với mỗi yêu cầu cho một số nguyên $K$, hãy xác định
\textbf{giá trị nhỏ thứ $K$} trong danh sách đánh giá nhóm.

\subsection*{Dữ liệu vào}

\begin{itemize}
    \item Dòng đầu tiên chứa số nguyên $T$ $(1 \le T \le 1000)$ --- số bộ dữ liệu.
    \item Với mỗi bộ dữ liệu:
    \begin{itemize}
        \item Dòng đầu chứa hai số nguyên $N$ và $Q$ $(3 \le N \le 3 \cdot 10^5, 1 \le Q \le 3 \cdot 10^5)$ --- số gia đình và số yêu cầu.
        \item Dòng tiếp theo chứa $N$ số nguyên: $A_1, A_2, \ldots, A_N$ $(-10^9 \le A_i \le 10^9)$ --- chỉ số sinh hoạt của các gia đình.
        \item $Q$ dòng tiếp theo, mỗi dòng chứa một số nguyên dương $K$ $(1 \le K \le \dfrac{N(N-1)(N-2)}{6})$.
    \end{itemize}
\end{itemize}

\subsection*{Dữ liệu ra}

Với mỗi yêu cầu, in ra \textbf{giá trị nhỏ thứ $K$}
trong danh sách đánh giá nhóm tương ứng.

\subsection*{Ví dụ}

\subsection*{Ví dụ}

\noindent
\begin{minipage}[t]{0.48\textwidth}
\textbf{Input}
\begin{verbatim}
2
6 2
3 3 3 3 3 3
18
1
4 4
2 4 2 1
1
2
3
4
\end{verbatim}
\end{minipage}
\hfill
\begin{minipage}[t]{0.48\textwidth}
\textbf{Output}
\begin{verbatim}
3
3
1
1
1
2
\end{verbatim}
\end{minipage}


\subsection*{Giải thích}

\textbf{Bộ dữ liệu 1.}
Mọi gia đình đều có cùng chỉ số sinh hoạt bằng $3$.
Do đó, với bất kỳ nhóm ba gia đình nào,
mức đánh giá thu được luôn là:
\[
3
\]

\textbf{Bộ dữ liệu 2.}
Các chỉ số sinh hoạt lần lượt là $[2,4,2,1]$.
Sau khi tổng hợp toàn bộ các nhóm ba gia đình
và sắp xếp các mức đánh giá thu được, ta có:
\[
[1,1,1,2]
\]

\newpage

\section*{Bài 4: Trạm Năng Lượng Tuyến Tính}

Trong một nền văn minh ngoài Trái Đất, người ta xây dựng một \emph{hành lang năng lượng} gồm $N$ mô-đun liên tiếp, được đánh số từ $1$ đến $N$.
Mỗi mô-đun chứa một lượng năng lượng không âm, được ghi nhận bởi một dãy số: $A_1, A_2, \dots, A_N$

Do đặc tính vật lý của không gian cong, hành lang này chỉ được coi là \emph{ổn định} nếu thỏa mãn tất cả các điều kiện sau:

\begin{itemize}
    \item Mô-đun đầu tiên phải \emph{đồng bộ} với mô-đun kế bên, hoặc hoàn toàn trống năng lượng:
    \[
    A_1 = A_2 \quad \text{hoặc} \quad A_1 = 0
    \]
    \item Mô-đun cuối cùng cũng phải tuân theo quy tắc tương tự:
    \[
    A_N = A_{N-1} \quad \text{hoặc} \quad A_N = 0
    \]
    \item Với mỗi mô-đun nằm giữa hành lang, năng lượng của nó phải được tạo ra từ hai mô-đun kề bên, \emph{hoặc} bị triệt tiêu hoàn toàn:
    \[
    A_i = A_{i-1} + A_{i+1}
    \quad \text{hoặc} \quad
    A_i = 0
    \qquad (1 < i < N)
    \]
\end{itemize}

Nói cách khác, bất kỳ mô-đun nào còn mang năng lượng đều phải được duy trì chính xác bằng tổng năng lượng của hai mô-đun lân cận.

Do sai lệch trong quá trình dịch chuyển không gian, hành lang ban đầu có thể ở trạng thái \emph{không ổn định}.
Hệ thống điều khiển cho phép thực hiện các thao tác hiệu chỉnh năng lượng như sau:

\begin{itemize}
    \item Tăng năng lượng của một mô-đun lên $1$ đơn vị, hoặc
    \item Giảm năng lượng của một mô-đun đi $1$ đơn vị,
\end{itemize}

với điều kiện rằng \textbf{không mô-đun nào được phép có năng lượng âm} tại bất kỳ thời điểm nào.

\textbf{Nhiệm vụ của bạn} là xác định \emph{số thao tác tối thiểu} cần thiết để đưa hành lang năng lượng về trạng thái ổn định.

Có thể chứng minh rằng, với mọi cấu hình ban đầu hợp lệ, luôn tồn tại ít nhất một phương án hiệu chỉnh.

\subsection*{Dữ liệu vào}

\begin{itemize}
    \item Dòng đầu tiên chứa số nguyên $T$ $(1 \le T \le 10^4)$ — số lượng kịch bản mô phỏng.
    \item Với mỗi kịch bản:
    \begin{itemize}
        \item Dòng đầu chứa số nguyên $N$ $(2 \le N \le 2 \cdot 10^5)$ — số mô-đun trong hành lang.
        \item Dòng tiếp theo chứa $N$ số nguyên không âm: $A_1, A_2, \dots, A_N \quad (0 \le A_i \le 10^9)$ — lượng năng lượng ban đầu của từng mô-đun.
    \end{itemize}
    \item Tổng $N$ của tất cả các kịch bản không vượt quá $2 \cdot 10^5$

\end{itemize}

\subsection*{Dữ liệu ra}

Với mỗi kịch bản, in ra một dòng chứa số thao tác tối thiểu cần thiết để đưa hành lang về trạng thái ổn định.

\subsection*{Ví dụ}

\noindent
\begin{minipage}[t]{0.48\textwidth}
\textbf{Input}
\begin{verbatim}
6
2
0 5
2
3 4
3
0 0 0
3
2 2 1
5
6 7 2 1 2
9
1000000000 1000000000 1000000000
1000000000 1000000000 1000000000
1000000000 1000000000 1000000000
\end{verbatim}
\end{minipage}
\hfill
\begin{minipage}[t]{0.48\textwidth}
\textbf{Output}
\begin{verbatim}
5
1
0
1
4
3000000000
\end{verbatim}
\end{minipage}

\subsection*{Giải thích}

\begin{itemize}
    \item \textbf{Kịch bản 1:}  
    Hành lang ban đầu là $[0, 5]$.  
    Việc triệt tiêu hoàn toàn năng lượng của mô-đun thứ hai đưa hệ thống về trạng thái ổn định.  
    Tổng số thao tác: $5$.

    \item \textbf{Kịch bản 2:}  
    Điều chỉnh mô-đun cuối để đồng bộ năng lượng với mô-đun đầu.  
    Tổng số thao tác: $1$.

    \item \textbf{Kịch bản 3:}  
    Tất cả mô-đun đều ở trạng thái trống năng lượng, hệ thống đã ổn định.  
    Không cần thao tác nào.
\end{itemize}

\bigskip
\begin{center}
\hdashrule{0.4\textwidth}{0.5pt}{2pt}\ \textbf{HẾT}\ \hdashrule{0.4\textwidth}{0.5pt}{2pt}

\noindent
\textit{Giám thị coi thi không giải thích gì thêm.}
\end{center}

